% TOP_PROBLEM: {"id": 3227, "title": "Circular coloring triangle-free subcubic planar graphs", "category_id": 3, "category": {"id": 3, "name": "graph_theory", "display_name": "Graph Theory", "description": "Problems involving graphs, networks, and their properties.", "slug": "graph-theory", "order_index": 3, "created_at": "2026-07-31T15:26:25.671Z"}, "status": "open", "problem_number": "OPG-401", "set_id": 12}
% FIRST_POSTED: 2026-10-01
% ENTRY_KIND: statement_only
% UnsolvedMath ID 3227; source catalogue code OPG-401
% !TeX program = lualatex
% Definition and sources only; this file is not an LLM solution attempt.
\documentclass[11pt,a4paper]{article}
\usepackage[margin=25mm]{geometry}
\usepackage{fontspec}
\setmainfont{lmroman10-regular.otf}[BoldFont=lmroman10-bold.otf,
 ItalicFont=lmroman10-italic.otf,BoldItalicFont=lmroman10-bolditalic.otf]
\usepackage{amsmath,amssymb,mathtools}
\usepackage{unicode-math}
\setmathfont{latinmodern-math.otf}
\AtBeginDocument{\let\setminus\smallsetminus}
\usepackage{xurl,hyperref}
\hypersetup{colorlinks=true,urlcolor=blue}
\setcounter{secnumdepth}{0}
\setlength{\parindent}{0pt}
\setlength{\parskip}{5pt}
\setlength{\emergencystretch}{3em}
\begin{document}
\section*{Circular coloring triangle-free subcubic planar graphs}\label{3227}
{\small UnsolvedMath ID 3227 · OPG-401 · Graph Theory · OpenGarden}
\subsection{Definitions and mathematical statement}
Let $G=(V,E)$ be a finite simple planar graph with maximum degree at most three and with no triangle. Planarity means existence of a drawing in the plane with disjoint edge interiors; triangle-free means there are no three pairwise adjacent vertices. No particular drawing is fixed and no lower bound on degree is assumed.

For integers $p\ge2q\ge2$, a circular $(p,q)$-coloring is a map $c:V\to\{0,\ldots,p-1\}$ satisfying $q\le|c(u)-c(v)|\le p-q$ on every edge. The circular chromatic number $\chi_c(G)$ is the infimum of $p/q$ over all such colorings. For an edgeless graph take $\chi_c(G)=1$.

The conjecture asserts
\[
                        \chi_c(G)\le\frac{20}{7}
\]
for every graph in this class. Equivalently, one asks for a map $c:V\to\mathbb Z_{20}$ such that each adjacent pair has cyclic color distance at least seven. Using representatives from $0$ through $19$, this condition is exactly $7\le|c(u)-c(v)|\le13$. This finite palette formulation removes any need to approximate the infimum in the target.

A vertex may receive the same label as any nonneighbor, and the coloring need not use all twenty labels. The assertion is stronger than the existence of an ordinary three-coloring: spacing around the circular palette imposes additional constraints. Disconnected graphs are included; the required coloring can be chosen component by component.
\subsection{Short English statement}
Can every triangle-free planar graph of maximum degree three be colored on a twenty-point circle with neighboring labels separated by at least seven steps?
\subsection{Sources}
\begin{itemize}
\item[S1] ULAM.AI and the UnsolvedMath Contributors, supplied problems.json, record 3227 (OPG-401), OpenGarden. Catalogue identity and source transcription; CC BY 4.0.
\url{https://huggingface.co/datasets/ulamai/UnsolvedMath}.
\item[S2] Open Problem Garden, Circular coloring triangle-free subcubic planar graphs. Original problem and discussion.
\url{http://www.openproblemgarden.org/op/circular_chromatic_number_of_triangle_free_planar_graphs}.
\end{itemize}
\end{document}
